% from ega4-2-fr.tex
\oldpage[IV-2]{131}

et $\mathfrak q$ est induit sur $A$ par un idéal premier de
$A_{\mathfrak p}$, minimal parmi ceux qui contiennent
$\mathfrak J'A_{\mathfrak p}$. Mais en vertu de l'hypothèse (ii) et de (0,
17.1.7), $\mathfrak J'A_{\mathfrak p}$ est premier dans $A_{\mathfrak p}$,
donc
\[
\mathfrak q=A\cap\mathfrak J'A_{\mathfrak p}
\]
est l'unique idéal premier non immergé associé à $A/\mathfrak J'$ et il est de
hauteur $n-1$, puisque $\mathfrak J'A_{\mathfrak p}$ est de hauteur $n-1$ et
$\mathfrak q$ de hauteur $\leq n-1$. En outre
$\mathfrak qA_{\mathfrak q}=\mathfrak J'A_{\mathfrak q}$ et comme
$A_{\mathfrak q}=(A_{\mathfrak p})_{\mathfrak qA_{\mathfrak p}}$, l'idéal
maximal $\mathfrak qA_{\mathfrak q}$ est engendré par $\mathfrak J'$. On voit
donc que la suite $(x_i)_{1\leq i\leq n-1}$ vérifie déjà les conditions (i) et
(ii) de l'énoncé. Montrons qu'elle vérifie aussi (iii). Considérons pour cela
l'anneau intègre $B=A/\mathfrak q$, et soit $t$ l'image canonique de $x_n$ dans
$B$. Il est immédiat (cf. (5.6.3, (iv))) que si $A$ vérifie l'hypothèse a)
(resp. b)), il en est de même de $B$; comme $x_n\notin\mathfrak q$, on a
$t\neq0$; montrons que $B$ et $t$ vérifient les hypothèses du lemme de
Hironaka. En effet, un idéal premier $\mathfrak n$ de $B$ minimal parmi ceux
contenant $t$ est l'image d'un idéal premier de $A$ minimal parmi ceux
contenant $\mathfrak q+x_nA$, et l'on a vu que $\mathfrak p$ est cet unique
idéal; comme
\[
B_{\mathfrak n}=A_{\mathfrak p}/\mathfrak qA_{\mathfrak p},\qquad
\mathfrak nB_{\mathfrak n}
=\mathfrak pA_{\mathfrak p}/\mathfrak qA_{\mathfrak p},
\]
le fait que $\mathfrak qA_{\mathfrak p}=\mathfrak J'A_{\mathfrak p}$ et que
$\mathfrak pA_{\mathfrak p}$ soit engendré par $\mathfrak J$ entraîne que
$\mathfrak nB_{\mathfrak n}$ est engendré par $t$. Enfin,
$B/\mathfrak n=A/\mathfrak p$ est intégralement clos. L'application de
(5.12.8) prouve donc que $B=A/\mathfrak q$ est intégralement clos et que
$\mathfrak p=\mathfrak q+x_nA$. Utilisons maintenant l'hypothèse de récurrence,
qui montre que $A/(x_1A+\cdots+x_iA)$ est intégralement clos et de dimension
$\dim(A)-i$ pour $0\leq i\leq n-1$ et que $\mathfrak q=\mathfrak J'$, d'où
$\mathfrak p=\mathfrak J'+x_nA=\mathfrak J$. On en conclut que
$A/\mathfrak J=A/\mathfrak p$ est intégralement clos; comme
$\dim(A_{\mathfrak p})=n$ et que
$\dim(A/\mathfrak p)=\dim(A)-\dim(A_{\mathfrak p})$ puisque $A$ est
biéquidimensionnel, cela achève la démonstration.

\subsection{Propriétés de permanence par passage à la limite inductive}

Dans tout ce numéro, $(A_\alpha,\varphi_{\beta\alpha})$ désigne un système
inductif d'anneaux dont l'ensemble d'indices est préordonné filtrant; on pose
$A=\varinjlim A_\alpha$ et on désigne par $\varphi_\alpha$ l'homomorphisme
canonique $A_\alpha\to A$.

\begin{proposition}[5.13.1]
\label{IV.2.s5.13.1-fr}
\begin{enumerate}
\item[(i)] Pour tout indice $\alpha$, soit $\mathfrak a_\alpha$ un idéal de
$A_\alpha$ tel que, pour $\alpha\leq\beta$, on ait
$\varphi_{\beta\alpha}(\mathfrak a_\alpha)\subset\mathfrak a_\beta$. Alors
les $\mathfrak a_\alpha$ forment un système inductif pour les restrictions des
$\varphi_{\beta\alpha}$, $\mathfrak a=\varinjlim\mathfrak a_\alpha$
s'identifie canoniquement à un idéal de $A$, $A/\mathfrak a$ à
$\varinjlim(A_\alpha/\mathfrak a_\alpha)$. Si de plus on a
$\mathfrak a_\alpha=\varphi_{\beta\alpha}^{-1}(\mathfrak a_\beta)$ pour
$\alpha\leq\beta$, alors
$\mathfrak a_\alpha=\varphi_\alpha^{-1}(\mathfrak a)$ pour tout $\alpha$.
\item[(ii)] Inversement, pour tout idéal $\mathfrak a$ de $A$, si l'on pose
$\mathfrak a_\alpha=\varphi_\alpha^{-1}(\mathfrak a)$ pour tout $\alpha$, les
$\mathfrak a_\alpha$ forment un système inductif tel que
$\mathfrak a=\varinjlim\mathfrak a_\alpha$.
\end{enumerate}
\end{proposition}

Il est clair que (ii) est un cas particulier de (i). Dans (i), le fait que
$\mathfrak a$ s'identifie canoniquement à un idéal de $A$ et $A/\mathfrak a$ à
$\varinjlim(A_\alpha/\mathfrak a_\alpha)$ provient de l'exactitude du foncteur
$\varinjlim$ dans la catégorie des groupes abéliens. En outre, $\mathfrak a$
est réunion de la famille croissante des $\varphi_\alpha(\mathfrak a_\alpha)$;
donc, si $x_\alpha\in A_\alpha$ est tel que
$\varphi_\alpha(x_\alpha)\in\mathfrak a$, il existe $\beta\geq\alpha$ tel que
$\varphi_\alpha(x_\alpha)=\varphi_\beta(y_\beta)$ pour un
$y_\beta\in\mathfrak a_\beta$, donc
$\varphi_\beta(\varphi_{\beta\alpha}(x_\alpha))=\varphi_\beta(y_\beta)$; mais
cela entraîne l'existence d'un $\gamma\geq\beta$ tel que
\[
\varphi_{\gamma\alpha}(x_\alpha)=\varphi_{\gamma\beta}(y_\beta),
\]
et par suite $\varphi_{\gamma\alpha}(x_\alpha)\in\mathfrak a_\gamma$, d'où
$x_\alpha\in\mathfrak a_\alpha$ si l'on a
$\mathfrak a_\alpha=\varphi_{\gamma\alpha}^{-1}(\mathfrak a_\gamma)$.

\begin{corollary}[5.13.2]
\label{IV.2.s5.13.2-fr}
Si $\mathfrak N_\alpha$ est le nilradical de $A_\alpha$, le nilradical de $A$
s'identifie à $\varinjlim\mathfrak N_\alpha$, et $A_{\mathrm{red}}$ à
$\varinjlim(A_\alpha)_{\mathrm{red}}$. En particulier, si les $A_\alpha$ sont
réduits, il en est de même de $A$.
\end{corollary}

